\documentclass[10pt,a4paper]{amsart}
\usepackage{amsmath,amssymb,amsthm,graphicx}
\usepackage{verbatim}
\newenvironment{reference}{\comment}{\endcomment}
%\newenvironment{reference}{}{}
\newenvironment{slogan}{\comment}{\endcomment}
\newenvironment{history}{\comment}{\endcomment}

% For commutative diagrams we use Xy-pic
\usepackage[all]{xy}

% We use 2cell for 2-commutative diagrams.
\xyoption{2cell}
% Theorem environments.
%
\theoremstyle{plain}
\newtheorem{theorem}[subsection]{Theorem}
\newtheorem{proposition}[subsection]{Proposition}
\newtheorem{lemma}[subsection]{Lemma}

\theoremstyle{definition}
\newtheorem{definition}[subsection]{Definition}
\newtheorem{example}[subsection]{Example}
\newtheorem{exercise}[subsection]{Exercise}
\newtheorem{situation}[subsection]{Situation}

\theoremstyle{remark}
\newtheorem{remark}[subsection]{Remark}
\newtheorem{remarks}[subsection]{Remarks}

\numberwithin{equation}{subsection}

% Macros
%
\def\lim{\mathop{\mathrm{lim}}\nolimits}
\def\colim{\mathop{\mathrm{colim}}\nolimits}
\def\Spec{\mathop{\mathrm{Spec}}}
\def\Hom{\mathop{\mathrm{Hom}}\nolimits}
\def\Ext{\mathop{\mathrm{Ext}}\nolimits}
\def\SheafHom{\mathop{\mathcal{H}\!\mathit{om}}\nolimits}
\def\SheafExt{\mathop{\mathcal{E}\!\mathit{xt}}\nolimits}
\def\Sch{\mathit{Sch}}
\def\Mor{\mathop{\mathrm{Mor}}\nolimits}
\def\Ob{\mathop{\mathrm{Ob}}\nolimits}
\def\Sh{\mathop{\mathit{Sh}}\nolimits}
\def\NL{\mathop{N\!L}\nolimits}
\def\CH{\mathop{\mathrm{CH}}\nolimits}
\def\proetale{{pro\text{-}\acute{e}tale}}
\def\etale{{\acute{e}tale}}
\def\QCoh{\mathit{QCoh}}
\def\Ker{\mathop{\mathrm{Ker}}}
\def\Im{\mathop{\mathrm{Im}}}
\def\Coker{\mathop{\mathrm{Coker}}}
\def\Coim{\mathop{\mathrm{Coim}}}

% Boxtimes
%
\DeclareMathSymbol{\boxtimes}{\mathbin}{AMSa}{"02}

%
% Macros for moduli stacks/spaces
%
\def\QCohstack{\mathcal{QC}\!\mathit{oh}}
\def\Cohstack{\mathcal{C}\!\mathit{oh}}
\def\Spacesstack{\mathcal{S}\!\mathit{paces}}
\def\Quotfunctor{\mathrm{Quot}}
\def\Hilbfunctor{\mathrm{Hilb}}
\def\Curvesstack{\mathcal{C}\!\mathit{urves}}
\def\Polarizedstack{\mathcal{P}\!\mathit{olarized}}
\def\Complexesstack{\mathcal{C}\!\mathit{omplexes}}
% \Pic is the operator that assigns to X its picard group, usage \Pic(X)
% \Picardstack_{X/B} denotes the Picard stack of X over B
% \Picardfunctor_{X/B} denotes the Picard functor of X over B
\def\Pic{\mathop{\mathrm{Pic}}\nolimits}
\def\Picardstack{\mathcal{P}\!\mathit{ic}}
\def\Picardfunctor{\mathrm{Pic}}
\def\Deformationcategory{\mathcal{D}\!\mathit{ef}}
\usepackage{cleveref}
\crefname{originaldefinition}{original definition}{original definitions}
\Crefname{originaldefinition}{Original definition}{Original definitions}
\makeatletter
\newlabel{sites-definition-morphism-sites}{{00X1}{source}{}{sourcepointer-1}{}}
\newlabel{sites-definition-morphism-sites@cref}{{[originaldefinition][0][]00X1}{[1][0][]source}{}{}{}}
\makeatother
\crefname{originallemma}{original lemma}{original lemmas}
\Crefname{originallemma}{Original lemma}{Original lemmas}
\makeatletter
\newlabel{sites-lemma-pullback-representable-sheaf}{{04D3}{source}{}{sourcepointer-2}{}}
\newlabel{sites-lemma-pullback-representable-sheaf@cref}{{[originallemma][0][]04D3}{[1][0][]source}{}{}{}}
\makeatother
\crefname{originalsection}{original section}{original sections}
\Crefname{originalsection}{Original section}{Original sections}
\makeatletter
\newlabel{sites-section-continuous-functors}{{00WU}{source}{}{sourcepointer-3}{}}
\newlabel{sites-section-continuous-functors@cref}{{[originalsection][0][]00WU}{[1][0][]source}{}{}{}}
\makeatother
\crefname{originalsection}{original section}{original sections}
\Crefname{originalsection}{Original section}{Original sections}
\makeatletter
\newlabel{stacks-cohomology-section-lisse-etale}{{0786}{source}{}{sourcepointer-4}{}}
\newlabel{stacks-cohomology-section-lisse-etale@cref}{{[originalsection][0][]0786}{[1][0][]source}{}{}{}}
\makeatother
\title{The lisse-etale site is not functorial}
\author{The Stacks Project Authors}
\date{stacks-project:examples.tex:local-snapshot-20261012; GNU FDL 1.2 or later}
\begin{document}
\maketitle
\noindent\textit{Editorial note.} All delivered mathematical text is the author's own, copied from the pinned Stacks Project chapter file examples.tex; no statement or proof step was rewritten, reordered or removed.\par\medskip
\noindent\textit{Editorial note.} The counted proof is the author's complete argument: Standard warning example for the lisse-etale topology on algebraic stacks.\par\medskip
\noindent\textit{Editorial note.} The label lemma-lisse-etale-not-functorial names the counted statement, and the proof environment immediately following it in the source is delivered verbatim as its proof.\par\medskip
\noindent\textit{Editorial note.} The counted statement is carried by the author's own lemma environment in the chapter source.\par\medskip
\noindent\textit{Allowed standard prerequisites (editorial).}\begin{itemize}
\item The lisse-etale site of a scheme X: the category of schemes smooth over X with the usual etale coverings (Cohomology of Stacks, Section stacks-cohomology-section-lisse-etale)
\item A continuous functor u between the underlying categories induces an adjoint pair of functors u-s and u-small on sheaf categories (Sites, Section sites-section-continuous-functors)
\item Sites, Definition sites-definition-morphism-of-sites: u defines a morphism of sites exactly when the pushforward u-small is left exact
\item Sites, Lemma sites-lemma-pullback-representable-sheaf: the pushforward of a representable sheaf under base change is the representable sheaf of the fibre product
\end{itemize}\medskip
\ifdefined\hypertarget\hypertarget{sourcepointer-1}{}\fi
\noindent\textit{Source reference.} Original definition 00X1: sources/stacks-chow/tags.txt, source lines 1203--1203. Prerequisite pointer invoked at 4822; it names a standard fact of a sibling chapter that the counted argument uses, and every step of the counted proof built on it is retained verbatim in the delivered body.\par\medskip
\ifdefined\hypertarget\hypertarget{sourcepointer-2}{}\fi
\noindent\textit{Source reference.} Original lemma 04D3: sources/stacks-chow/tags.txt, source lines 5713--5713. Prerequisite pointer invoked at 4825; it names a standard fact of a sibling chapter that the counted argument uses, and every step of the counted proof built on it is retained verbatim in the delivered body.\par\medskip
\ifdefined\hypertarget\hypertarget{sourcepointer-3}{}\fi
\noindent\textit{Source reference.} Original section 00WU: sources/stacks-chow/tags.txt, source lines 1196--1196. Prerequisite pointer invoked at 4820; it names a standard fact of a sibling chapter that the counted argument uses, and every step of the counted proof built on it is retained verbatim in the delivered body.\par\medskip
\ifdefined\hypertarget\hypertarget{sourcepointer-4}{}\fi
\noindent\textit{Source reference.} Original section 0786: sources/stacks-chow/tags.txt, source lines 9226--9226. Prerequisite pointer invoked at 4801; it names a standard fact of a sibling chapter that the counted argument uses, and every step of the counted proof built on it is retained verbatim in the delivered body.\par\medskip
\label{section-lisse-etale-not-functorial}

\noindent
The {\it lisse-\'etale} site
$X_{lisse,\etale}$ of $X$ is the category of schemes smooth over $X$
endowed with (usual) \'etale coverings, see
Cohomology of Stacks, Section \ref{stacks-cohomology-section-lisse-etale}.
Let $f : X  \to Y$ be a morphism of schemes.
There is a functor
$$
u : Y_{lisse,\etale} \longrightarrow X_{lisse,\etale},\quad
V/Y \longmapsto V \times_Y X
$$
which is continuous. Hence we obtain an adjoint pair of functors
$$
u^s :
\Sh(X_{lisse,\etale})
\longrightarrow
\Sh(Y_{lisse,\etale}),
\quad
u_s :
\Sh(Y_{lisse,\etale})
\longrightarrow
\Sh(X_{lisse,\etale}),
$$
see Sites, Section \ref{sites-section-continuous-functors}.
We claim that, in general, $u$ does {\bf not} define a morphism of sites, see
Sites, Definition \ref{sites-definition-morphism-sites}.
In other words, we claim that $u_s$ is not left exact in general. Note that
representable presheaves are sheaves on lisse-\'etale sites. Hence, by
Sites, Lemma \ref{sites-lemma-pullback-representable-sheaf}
we see that $u_sh_V = h_{V \times_Y X}$. Now consider two morphisms
$$
\xymatrix{
V_1 \ar[rd] \ar@<1ex>[rr]^a \ar@<-1ex>[rr]_b & & V_2 \ar[ld] \\
& Y
}
$$
of schemes $V_1, V_2$ smooth over $Y$. Now if $u_s$ is left exact, then
we would have
$$
u_s \text{Equalizer}(h_a, h_b : h_{V_1} \to h_{V_2})
=
\text{Equalizer}(h_{a \times 1}, h_{b \times 1} :
h_{V_1 \times_Y X} \to h_{V_2 \times_Y X})
$$
We will take the morphisms $a, b : V_1 \to V_2$ such that there exists
no morphism from a scheme smooth over $Y$ into $(a = b) \subset V_1$, i.e.,
such that the left hand side is the empty sheaf, but such that after
base change to $X$ the equalizer is nonempty and smooth over $X$.
A silly example is to take $X = \Spec(\mathbf{F}_p)$,
$Y = \Spec(\mathbf{Z})$ and $V_1 = V_2 = \mathbf{A}^1_\mathbf{Z}$
with morphisms $a(x) = x$ and $b(x) = x + p$. Note that the equalizer
of $a$ and $b$ is the fibre of $\mathbf{A}^1_\mathbf{Z}$ over $(p)$.


\begin{lemma}
\label{lemma-lisse-etale-not-functorial}
The lisse-\'etale site is not functorial, even for morphisms of schemes.
\end{lemma}

\begin{proof}
See discussion above.
\end{proof}

\end{document}
